\chapter{Tutorial 2: from a Gaussian anharmonic calculation}\label{ch:tut-gaussian}

This tutorial starts from the output of a Gaussian
\texttt{freq=(anharmonic,...)} calculation, here for
1,2,3,6,7,8-hexahydropyrene (90 modes, B3LYP-D3/6-31G). The Gaussian input
and the compressed output are in \file{examples/hexahydropyrene}; the
same calculation is published in the PAH database of INAF--OACa,
\url{https://cosmicpah-qcals.oa-cagliari.inaf.it}. The route section is
\begin{lstlisting}
#p geom=check freq=(ReadFC,anharm,PrintDerivatives) b3lyp/gen/auto
   int=(grid=superfine) cphf=(grid=superfine) densityfit
   empiricaldispersion=gd3 iop(7/33=1,99/5=0) scf=xqc
\end{lstlisting}
The essential keywords are \texttt{freq=(anharm,PrintDerivatives)}, which
prints the cubic and quartic force constants and the dipole derivatives
in normal coordinates. \texttt{ReadFC} and \texttt{geom=check} reuse a
previous optimisation and frequency calculation.

\section{Converting the Gaussian output}

\begin{lstlisting}
xz -dk 1_2_3_6_7_8-hexahydropyrene_distort1b.log.xz
python3 tools/gaussian2ac.py 1_2_3_6_7_8-hexahydropyrene_distort1b.log
\end{lstlisting}
The script writes \file{freqsfile}, \file{phi3file}, \file{phi4file},
\file{dmu2file} and, when the log contains them, the semidiagonal third
derivatives of the dipole (\file{dmu3file}, used with \opt{dipole3}). The
modes are renumbered by increasing wavenumber.

\section{Comparing with Gaussian}

\file{dunhamparams} uses the same convention as Gaussian's
``Total Anharmonic X Matrix'': each pair $i\le j$ is counted once.
Gaussian's matrix includes the Coriolis contributions, which \ac{}
neglects. Once these are subtracted, the two agree to the printed
precision, except for the entries that Gaussian deperturbs for its Fermi
resonances (129 of 4095 for this molecule). The zero-point energy agrees
to 0.03~cm$^{-1}$.
% TODO: script for this comparison; tables of fundamentals vs GVPT2.
